\subsection*{可解李代数的例子}

I. 设 g 是 K 上的二维向量空间，\((e_{1}, e_{2})\) 是 g 的一组基。g 上存在唯一的交错双线性乘法 \((x, y) \mapsto [x, y]\)，使得 \([e_{1}, e_{2}] = e_{2}\)。容易验证，这样 g 就成为可解李代数。现在设 h 是 K 上二维的非交换李代数。我们证明 h 同构于 g。设 \((f_{1}, f_{2})\) 是 h 的一组基。元素 \([f_{1}, f_{2}]\) 非零（否则 h 将交换），因而它生成 h 的一维子空间 t。于是 \([h, h] = t\)。设 \((e_{1}^{\prime}, e_{2}^{\prime})\) 是 h 的一组基，且 \(e_{2}^{\prime} \in t\)。则 \([e_{1}^{\prime}, e_{2}^{\prime}] = \lambda e_{2}^{\prime}\)，其中 \(\lambda \neq 0\)。将 \(e_{1}^{\prime}\) 替换为 \(\lambda^{-1}e_{1}\)，可设 \(\lambda = 1\)，从而得到所述断言。

\begin{enumerate}
\def\labelenumi{\Roman{enumi}.}
\setcounter{enumi}{1}
\tightlist
\item
  §1 的公式 (5) 证明了 \(\mathcal{D}t(n, \mathbf{K}) = \mathfrak{n}(n, \mathbf{K})\)。由于 \(\mathfrak{n}(n, \mathbf{K})\) 是幂零的，因而可解，\(t(n, \mathbf{K})\) 是可解的。因此 \(\operatorname{st}(n, \mathbf{K})\) 是可解的。特别地，\(\operatorname{st}(2, \mathbf{K})\) 同构于例 I 中的代数。
\end{enumerate}

